\section{Literature Survey}

We reviewed seven studies across three modality groups: MRI, EEG and voice (Table~\ref{tab:lit}). MRI
studies report very high accuracies but rely on expensive imaging, are often evaluated on 2D slices, and
rarely validate across datasets. EEG-based approaches remain underexplored by comparison, and none of the
reviewed studies validates on an independent dataset.

\begin{table}[H]
\centering
\caption{Literature survey: seven studies across three modality groups}
\label{tab:lit}
\footnotesize
\begin{tabularx}{\textwidth}{@{}>{\raggedright\arraybackslash}p{2.6cm} >{\raggedright\arraybackslash}p{2.6cm}
  >{\raggedright\arraybackslash}X >{\raggedright\arraybackslash}p{2.0cm} >{\raggedright\arraybackslash}X@{}}
\toprule
Study & Data & Method & Reported result & Limitation\\
\midrule
\multicolumn{5}{@{}l}{\textit{MRI-based approaches}}\\
\citet{elassy2024cnn} & ADNI & Dual CNN branches with different filter sizes, concatenated; 3/4/5-class & 99.43\%, 99.57\%, 99.13\% & Possible overfitting; no cross-dataset validation; limited explainability\\
\citet{odimayo2024structure} & ADNI + OASIS & Structure-focused CNN, gamma correction, midsagittal and parasagittal planes & 97--98\% & Specific MRI planes only; high computational cost; limited generalisation testing\\
\citet{mousavi2025classifying} & ADNI-derived MRI (6,735 images) & Transfer learning (VGG16, VGG19, Xception, InceptionResNetV2), majority-vote ensemble & 98.52\% & Derived dataset may lack clinical diversity; mostly 2D slice-based\\
\midrule
\multicolumn{5}{@{}l}{\textit{EEG-based approaches}}\\
\citet{kim2023caueeg} & CAUEEG & CEEDNet: end-to-end 1D/2D CNN and transformer family on raw EEG with age & 70.2\% single, 74.7\% ensemble (three-class, per recording) & Single dataset; large models (11--87M parameters); limited interpretability\\
\citet{miltiadous2023dataset} & AD, FTD and healthy EEG & Dataset publication & Not a model paper & No benchmark model; limited size\\
\citet{jungrungrueang2025translational} & Resting-state EEG & Dynamic EEG connectivity features with CNN classification & 93.6\% (multi-class) & EEG only; no MRI integration\\
\midrule
\multicolumn{5}{@{}l}{\textit{Other modalities}}\\
\citet{xue2021voice} & Framingham Heart Study voice recordings & CNN + LSTM on audio, minimal feature engineering & Binary classification & No MRI comparison; restricted clinical data\\
\bottomrule
\end{tabularx}\\[2pt]
\parbox{\textwidth}{\scriptsize ADNI: Alzheimer's Disease Neuroimaging Initiative; OASIS: Open Access Series of
Imaging Studies; CNN: convolutional neural network; LSTM: long short-term memory; AD: Alzheimer's disease; FTD:
frontotemporal dementia.}
\end{table}

\textbf{EEG markers of cognitive decline.} Spectral slowing (theta/alpha and delta/alpha ratios), lower alpha
peak frequency, reduced alpha reactivity to eye opening and loss of inter-regional coherence are consistently
reported in Alzheimer's disease and, more weakly, in MCI \citep{jeong2004eeg,dauwels2010diagnosis,babiloni2020measures}.

\textbf{Machine learning on CAUEEG.} \citet{kim2023caueeg} released CAUEEG together with CEEDNet, trained on
random 10\,s crops with age as an input. CEFA-Net \citep{cefanet2026} combines a three-scale CNN with a branch
of clinical EEG ratios and a focal loss, reporting 81.0\% with per-segment evaluation and 73.8\% on a
patient-disjoint test. A microstate-based support vector machine reports 89\% \citep{hasan2025msaf}. These
numbers use different evaluation units and splits and are not directly comparable. \citet{geneeg2025}
contrasted CAUEEG with GENEEG, a general-practice cohort, and found poor transfer between the two and strongly
overlapping MCI and control feature distributions. Compact architectures such as EEGNet
\citep{lawhern2018eegnet} are common baselines for small clinical datasets.

\textbf{Normative modelling and shortcut learning.} Normative models describe each subject as a deviation from
a reference population conditioned on covariates such as age \citep{marquand2016normative}, which separates
normal ageing from pathology. Classifiers frequently exploit acquisition artifacts instead of the target
signal \citep{geirhos2020shortcut}; in medical imaging, scanner or protocol differences between cohorts are a
well-known source \citep{roberts2021pitfalls,varoquaux2022machine}.

\textbf{Research gap.} No reviewed work combines (i) a lightweight EEG model that uses recording-state
information, (ii) a comparison with strong baselines under an identical, patient-disjoint protocol, and
(iii) validation on independent datasets. This work addresses all three.
